\documentclass[12pt]{article}

% Core mathematics
\usepackage{amsmath,amssymb,amsthm}
\usepackage{mathtools}

% Diagrams
\usepackage{tikz-cd}
\usepackage{tikz}

% References
\usepackage{hyperref}
\usepackage[nameinlink,noabbrev]{cleveref}

% Tables and layout
\usepackage{graphicx}
\usepackage{booktabs}
\usepackage{longtable}
\usepackage{array}
\usepackage[margin=1in]{geometry}
\usepackage{microtype}
\usepackage{enumitem}

% GrokRxiv sidebar
\usepackage{everypage}
\usepackage{xcolor}

\hypersetup{
  colorlinks=true,
  linkcolor=blue!55!black,
  citecolor=green!40!black,
  urlcolor=blue!60!black,
  pdftitle={The Invertible Condensed Phase Spectrum Program},
  pdfauthor={Matthew Long and the YonedaAI Collaboration}
}

\newtheorem{theorem}{Theorem}[section]
\newtheorem{proposition}[theorem]{Proposition}
\newtheorem{lemma}[theorem]{Lemma}
\newtheorem{corollary}[theorem]{Corollary}
\theoremstyle{definition}
\newtheorem{definition}[theorem]{Definition}
\newtheorem{example}[theorem]{Example}
\newtheorem{counterexample}[theorem]{Counterexample}
\theoremstyle{remark}
\newtheorem{remark}[theorem]{Remark}

\newcommand{\A}{\mathcal A}
\newcommand{\B}{\mathcal B}
\newcommand{\C}{\mathcal C}
\newcommand{\F}{\mathcal F}
\newcommand{\Ham}{\mathfrak{Ham}}
\newcommand{\Gap}{\mathfrak{Gap}}
\newcommand{\Phase}{\mathfrak{Phase}}
\newcommand{\Pic}{\operatorname{Pic}}
\newcommand{\Core}{\operatorname{Core}}
\newcommand{\Map}{\operatorname{Map}}
\newcommand{\Spec}{\operatorname{Spec}}
\newcommand{\State}{\operatorname{State}}
\newcommand{\Cont}{\operatorname{Cont}}
\newcommand{\Th}{\operatorname{Th}}
\newcommand{\im}{\operatorname{im}}
\newcommand{\IP}{\mathit{IP}}
\newcommand{\FH}{\mathit{FH}}
\newcommand{\Cond}{\mathit{Cond}}
\newcommand{\Kop}{K_{\mathrm{op}}}
\newcommand{\Kalg}{K_{\mathrm{alg}}}
\newcommand{\EST}{\textnormal{\textsc{Established}}}
\newcommand{\PROP}{\textnormal{\textsc{Proposed}}}
\newcommand{\CONJ}{\textnormal{\textsc{Conjectural}}}
\newcommand{\OPEN}{\textnormal{\textsc{Open}}}
\newcommand{\OBST}{\textnormal{\textsc{Obstructed}}}

\definecolor{grokgray}{RGB}{110,110,110}
\AddEverypageHook{%
  \ifnum\value{page}=1
    \begin{tikzpicture}[remember picture,overlay]
      \node[
        rotate=90,
        anchor=south,
        font=\Large\sffamily\bfseries\color{grokgray},
        inner sep=0pt
      ] at ([xshift=38pt,yshift=0.52\paperheight]current page.south west)
      {GrokRxiv:2026.07.invertible-condensed-phase-spectrum\quad
       [\,math-ph\,]\quad
       14 Jul 2026};
    \end{tikzpicture}
  \fi
}

\title{The Invertible Condensed Phase Spectrum Program:\\
A Conditional Synthesis of Locality, Gaps, Phases, and Realization}
\author{Matthew Long \\
\textit{The YonedaAI Collaboration} \\
\textit{YonedaAI Research Collective} \\
Chicago, IL \\
\texttt{matthew@yonedaai.com} $\cdot$ \url{https://yonedaai.com}}
\date{14 July 2026}

\begin{document}

\maketitle

\begin{abstract}
This paper asks whether local quantum interactions can be organized into an
invertible condensed phase spectrum without erasing the analytic hypotheses
that make a lattice model physical.  The proposed construction has five
stages:
\[
\boxed{
\begin{array}{c}
\text{local quantum interactions}\\
\downarrow\\
\text{condensed moduli stack of Hamiltonians}\\
\downarrow\\
\text{uniformly gapped propositional substack}\\
\downarrow\\
\text{stabilized phase }\infty\text{-groupoid}\\
\downarrow\\
\text{invertible condensed phase spectrum}.
\end{array}}
\]
The sequence is a research program, not a chain of established equivalences.
Its analytic input is solid in several controlled settings.  Common
interaction norms give Lieb-Robinson bounds.  Family-uniform thermodynamic
dynamics additionally require a common interaction tail that vanishes at
large distance.
Continuous parameter families in a fixed C*-algebra preserve positivity and
the C*-norm.  Uniform finite-volume or GNS gaps support quasi-adiabatic
spectral flow and phase stability.  Functional calculus maps controlled free
families to operator K-theory.  Low-dimensional Kitaev, BDI, and AKLT models
provide explicit realizability witnesses.

The categorical assembly is less complete.  A quantitative gap witness is
data, so its projection to Hamiltonians is a fibration and not a subobject.
The qualitative gapped locus is instead the propositional image of that
projection.  Localization by controlled gapped paths, stabilization by
atomic ancillas, stacking, and passage to the invertible sector should then
produce a grouplike symmetric monoidal condensed anima.  A connective
spectrum follows if these constructions satisfy descent and the required
coherences.  Those hypotheses remain proposed or conjectural.

We keep the resulting object distinct from operator K-theory, Kubota's
microscopic Omega-spectrum, and the Freed-Hopkins bordism-based field-theory
classification.  No identification is made without a comparison theorem.
For a family on $B$ with gapless locus $\Sigma$ and
$U=B\setminus\Sigma$, a class $\nu\in E^q(U)$ has the established relative
boundary $\partial\nu\in E^{q+1}(B,U)$.  Its interpretation as a physical
transition charge is conditional on a microscopic comparison map.  Excision
and an $E$-oriented normal bundle of codimension $c$ identify it with a class
in $E^{q+1-c}(\Sigma)$; without orientation the target is twisted.

The result is a precise conditional assembly theorem, a claim-status ledger,
and an executable Haskell validator that rejects the principal false
inferences.  It does not claim a completed condensed spectrum or a universal
classification of topological matter.
\end{abstract}

\tableofcontents

\section{The research question}

Topological phases are commonly described by stable labels.  Depending on
the problem, the label may be a Chern number, a K-class, a bordism invariant,
or a homotopy group of a space of gapped systems.  A label is useful because
it survives a controlled deformation.  It is dangerous when the deformation
rules and the comparison map that produced it are left implicit.

The condensed-mathematics perspective begins one step earlier
\cite{Scholze2026}.  It asks for
a moduli object whose test points are parameterized families of local
Hamiltonians.  Profinite tests are especially natural for disorder hulls,
inverse limits, and totally disconnected parameter spaces.  Descent should
glue compatible local descriptions, while higher morphisms should remember
paths, homotopies of paths, pumps, and defects.

The program studied here can be stated in one sentence.

\begin{quote}
Construct a condensed moduli object of uniformly local Hamiltonian families,
cut out its uniformly gapped qualitative locus, impose the physical phase
relation and stabilization, and spectrify the stacking-invertible sector.
\end{quote}

Every verb in that sentence hides a theorem obligation.  ``Construct'' asks
for a site, objects, morphisms, and descent.  ``Cut out'' asks whether a gap
condition is stable under restriction and gluing.  ``Impose'' asks for a
calculus of controlled paths and quasi-local equivalences.  ``Stabilize''
asks which ancillas count as trivial.  ``Spectrify'' asks for coherent
symmetric monoidal structure and deloopings.  The paper keeps those
obligations visible.

\subsection{Five stages, not five equivalences}

Write the proposed sequence as
\begin{equation}
\mathsf{Int}^{\mathrm{ub}}_{\F}
\longrightarrow
\Ham_{\F}
\longleftarrow
\Gap^{\mathrm{prop}}_{\F}
\longrightarrow
\Phase^{\mathrm{st}}_{\F}
\longrightarrow
\IP^{\mathrm{cond}}_{\F}.
\label{eq:five-stage-symbolic}
\end{equation}
The reverse arrow at the gapped stage is intentional: a gapped subobject,
if it exists, includes into the Hamiltonian moduli object.  A quantitative
witness object maps to both.  The prose sequence records a flow of
construction; the categorical diagram records the variance correctly.

No arrow in \cref{eq:five-stage-symbolic} is declared an equivalence.  The
first packages analytic data.  The second takes a propositional image.  The
third localizes and stabilizes.  The fourth passes to stacking-invertible
objects and then, conditionally, to a spectrum.

\subsection{Spatial and spacetime grading}

We use
\begin{equation}
d=\text{spatial dimension},
\qquad
n=d+1=\text{spacetime dimension}.
\label{eq:dimension-convention}
\end{equation}
Microscopic interactions and Kubota's construction are graded by $d$.
Field-theory tangential structures are graded by $n$.  Any comparison must
display the shift.  Suppressing it can make two formulas look identical even
when they classify different dimensional objects.

\subsection{Scope}

The target is the invertible sector.  Intrinsic topological order with
noninvertible excitations needs fusion, braiding, and defect categories that
cannot be recovered from a spectrum of units.  Crystalline symmetry,
fermionic grading, and disorder may be included only after their actions are
specified in the microscopic category.  This paper describes interfaces for
that work rather than claiming a complete classification.

\section{A status grammar for a mixed theorem program}

The source papers combine imported theorems, finite calculations,
definitions, and proposed categorical constructions.  A single word such as
``result'' cannot distinguish them.  We use five statuses.

\begin{definition}[Claim status]
For a statement made under explicit hypotheses:
\begin{itemize}[leftmargin=2.1em]
  \item \EST{} means a theorem or complete calculation is available in the
  stated setting.
  \item \PROP{} means a definition or construction is specified, but its
  full existence, functoriality, or usefulness has not been proved.
  \item \CONJ{} means a mathematically definite statement is expected to be
  true and lacks a proof.
  \item \OPEN{} means the paper asserts neither a construction nor a
  resolution.
  \item \OBST{} means the stated implication is false or is ruled out under
  the stated rules.
\end{itemize}
\end{definition}

The difference between \OPEN{} and \OBST{} matters.  A missing lattice model
is not an obstruction.  An anomaly theorem or an explicit interacting
counterexample may be one.  The difference between \PROP{} and \CONJ{} also
matters.  The notation $\IP^{\mathrm{cond}}$ is a proposed target.  A claim
that it agrees with another spectrum on a specified overlap is conjectural
only after a comparison map and its expected domain have been stated.

\subsection{An evidence order}

Statuses do not form a simple scale from weak to strong.  An obstructed
implication can be known more firmly than a proposed construction.  For
validation purposes it is better to record provenance separately:
\begin{equation}
\text{claim}\quad+\quad\text{status}\quad+\quad
\text{hypotheses}\quad+\quad\text{provider}.
\label{eq:claim-record}
\end{equation}
A provider may be a theorem, an explicit calculation, a cited construction,
or a named open comparison problem.  The Haskell code mirrors
\cref{eq:claim-record}; it never tries to infer truth from a status label.

\subsection{Nonclaims fixed at the outset}

The following statements are excluded.
\begin{enumerate}[leftmargin=2.2em]
  \item A profinite parameter space does not imply a uniform locality bound.
  \item Pointwise finite-volume gaps do not imply a thermodynamic gap.
  \item A ring with involution is not automatically a C*-algebra.
  \item A K-class does not reconstruct a Hamiltonian.
  \item A low-energy field theory is not automatically equivalent to its
  microscopic lattice model.
  \item Isomorphic homotopy groups do not identify spectra.
  \item An abstract bordism class does not provide a microscopic realization.
  \item A relative cohomology class is not automatically a measurable
  transition charge.
\end{enumerate}

These exclusions are not philosophical caution.  Each one corresponds to a
missing hypothesis or a known loss of information.

\section{Stage one: local quantum interactions}

Let $\Gamma$ be a countable metric lattice of bounded geometry.  At each site
$x$ choose a finite-dimensional Hilbert space $\mathcal H_x$, or a finite
dimensional graded local algebra for fermions.  The quasi-local algebra
$\A_\Gamma$ is the norm completion of the directed union of finite-region
matrix algebras.

An interaction is a map
\begin{equation}
\Phi:\{X\Subset\Gamma\}\longrightarrow(\A_X)_{\mathrm{sa}}.
\end{equation}
Choose a decay function $F:[0,\infty)\to(0,\infty)$ with the summability and
convolution properties used in Lieb-Robinson theory.  One standard norm has
the form
\begin{equation}
\|\Phi\|_F
=
\sup_{x,y\in\Gamma}
\frac{1}{F(d(x,y))}
\sum_{X\ni x,y}\|\Phi(X)\|.
\label{eq:f-norm}
\end{equation}
Other weighted norms are possible.  What matters in a family is that one
norm convention and one bound work for every parameter.

\begin{definition}[Uniformly bounded local family]
For a compact or profinite parameter object $S$, a family
$\Phi=(\Phi_s)_{s\in S}$ is uniformly $F$-local when
\begin{equation}
\sup_{s\in S}\|\Phi_s\|_F<\infty
\label{eq:uniform-family-locality}
\end{equation}
and its local coefficients $s\mapsto\Phi_s(X)$ are norm-continuous for every
finite $X$.  When infinite tails are admitted, the family also carries an
explicit tail bound uniform in $s$.
\end{definition}

Local coefficient continuity and \cref{eq:uniform-family-locality} are
independent conditions.  A family can vary continuously on every finite
region while its interaction range escapes to infinity.  Conversely, a
uniformly finite-range assignment can jump discontinuously in its
coefficients.

\subsection{Established locality output}

\begin{theorem}[Uniform Lieb-Robinson control]
\label{thm:uniform-lr}
\EST{}  Suppose a family of interactions has a common finite $F$-norm on a
bounded-geometry lattice.  Then the
finite-volume dynamics satisfy a Lieb-Robinson estimate with constants that
can be chosen uniformly over the family.  For local observables $A$ and $B$
with separated supports, one obtains a bound of the form
\begin{equation}
\|[\tau^{\Lambda,s}_t(A),B]\|
\le
C\|A\|\|B\|
\bigl(e^{v|t|}-1\bigr)
\sum_{x\in\operatorname{supp}A}
\sum_{y\in\operatorname{supp}B}F(d(x,y)),
\label{eq:lr-bound}
\end{equation}
where $C$ and $v$ do not depend on $s$ or on the finite volume $\Lambda$
\cite{LiebRobinson1972,NachtergaeleSims2010}.  For the thermodynamic limit of
the full family, require separately a common tail function $T_F(R)$ with
$T_F(R)\to0$ as $R\to\infty$.  Finite-region $F$-norm continuity does not
imply this uniform tail condition.
\end{theorem}

This theorem is analytic input, not a consequence of condensed descent.  A
test object organizes the parameter dependence; it does not manufacture the
common constant in \cref{eq:uniform-family-locality}.

\subsection{Restriction and finite descent}

For a map $f:T\to S$, restriction sends $\Phi$ to $f^*\Phi$.  The same
uniform locality bound remains valid.  For a finite jointly surjective cover
$S_i\to S$ of profinite spaces, compatible local coefficient functions glue
because the disjoint union map is a quotient map.  A common numerical bound
also glues if it has been supplied on the entire finite cover, for example by
taking the maximum of finitely many bounds.

This gives a controlled set-valued sheaf of interaction objects on the
chosen finite-cover site.  It does not yet give the higher moduli stack used
later.  Morphisms, automorphisms, and higher coherences still have to be
defined.

\begin{counterexample}[Moving weak defect]
Let $S=\mathbb N\cup\{\infty\}$ be the one-point compactification.  A family
may place a weak local perturbation farther from the origin as the parameter
approaches $\infty$.  Every fixed local coefficient then converges, and the
interactions may share a locality norm.  A low-energy mode tied to the moving
defect can nevertheless make the spectral gap tend to zero.  Local
continuity and uniform locality therefore do not imply a uniform gap.
\end{counterexample}

\section{Stage two: a condensed Hamiltonian moduli object}

The interaction sheaf records objects.  A moduli stack must also record
isomorphisms and higher families.  Its definition depends on which changes
of microscopic presentation count as equivalences.

\subsection{Objectwise C*-algebraic control}

For a compact Hausdorff parameter space $S$ and a C*-algebra $\A$, the
continuous family algebra
\begin{equation}
C(S,\A)
\end{equation}
is a C*-algebra with pointwise operations, involution, and norm
\begin{equation}
\|a\|=\sup_{s\in S}\|a(s)\|_{\A}.
\end{equation}
Positivity is fiberwise:
\begin{equation}
a\ge0
\quad\Longleftrightarrow\quad
a(s)\ge0\text{ for every }s\in S.
\label{eq:fiberwise-positive}
\end{equation}
The C*-identity follows exactly:
\begin{equation}
\|a^*a\|
=
\sup_s\|a(s)^*a(s)\|
=
\sup_s\|a(s)\|^2
=
\|a\|^2.
\label{eq:family-cstar}
\end{equation}

These statements solve the positivity and norm problem once the fixed
C*-algebra $\A$ has been chosen.  They do not construct that norm from bare
algebraic multiplication and involution.

\begin{counterexample}[A star algebra can have several C*-completions]
A dense involutive convolution algebra associated with a nonamenable group
can admit distinct full and reduced C*-norms.  Its multiplication and star
operation do not select one completion.  Any Hamiltonian moduli problem that
uses spectra, positivity, or states must carry the completion as part of its
input.
\end{counterexample}

\subsection{Finite-cover C*-descent}

If $S_i\to S$ is a finite jointly surjective family of profinite spaces, then
\begin{equation}
C(S,\A)
\cong
\operatorname{Eq}\left(
\prod_i C(S_i,\A)
\rightrightarrows
\prod_{i,j}C(S_i\times_S S_j,\A)
\right)
\label{eq:cstar-descent}
\end{equation}
as C*-algebras.  The norm on a glued section is forced by the fiberwise
supremum.  Positivity and self-adjointness descend because they are tested
fiberwise.

Finite clopen partitions give dense finite approximations to continuous
functions on a profinite space.  They do not imply that every continuous
family factors exactly through a finite quotient.  That stronger assertion
would discard genuine continuous variation.

\subsection{States are not pointwise state-valued maps}

A state on $C(S,\A)$ contains a probability measure on $S$ together with
conditional state data.  In general,
\begin{equation}
\State(C(S,\A))
\ne
C(S,\State(\A)).
\label{eq:state-warning}
\end{equation}
The right side chooses one normalized state at every point.  The left side
can average over the parameter space.  This distinction is relevant when a
ground-state sector is part of a gap witness.

\subsection{The proposed moduli stack}

Let $\mathsf{CondTest}$ be a chosen site of condensed test objects.  Its
precise size, topology, and hypercompletion convention must be fixed.  For
each $S$, define a tentative category $\C_{\Ham}(S)$ whose objects include:
\begin{enumerate}[leftmargin=2.2em]
  \item a lattice or relaxed-lattice datum of spatial dimension $d$;
  \item a continuous field of local finite-dimensional algebras;
  \item a uniformly $F$-local self-adjoint interaction;
  \item symmetry and grading data;
  \item boundary and ground-state conventions when required.
\end{enumerate}
Morphisms should record controlled changes of presentation, such as local
unitary identifications, lattice refinements, or paths of interactions.  The
choice determines the eventual phase relation.

\begin{definition}[Condensed Hamiltonian moduli, proposed]
\PROP{}  The object $\Ham_{d,\F}^G$ is the hypercomplete sheafification of
the anima-valued presheaf
\begin{equation}
S\longmapsto\Core(\C_{\Ham}(S)),
\label{eq:ham-stack-proposed}
\end{equation}
provided restriction preserves the analytic bounds and the chosen
morphisms satisfy descent.
\end{definition}

The word ``provided'' carries real content.  Objectwise finite-cover descent
does not automatically prove descent for path spaces, automorphism groups,
spectral-flow providers, or variable lattices.  The existence of
\cref{eq:ham-stack-proposed} with all desired structures is \PROP{}.

\section{Stage three: the uniformly gapped propositional substack}

A thermodynamic gap is quantitative.  A gapped locus is qualitative.  The
two are related by a projection, but they are not the same object.

\subsection{Finite-volume and GNS witnesses}

Fix a family of finite volumes $\Lambda_L$ and a prescribed low-energy
sector of dimension $m_L$.  A uniform finite-volume witness contains a
number $\gamma>0$ and a threshold $L_0$ such that
\begin{equation}
E_{m_L+1}(H_{\Lambda_L,s})-E_{m_L}(H_{\Lambda_L,s})
\ge\gamma
\label{eq:finite-volume-family-gap}
\end{equation}
for every $L\ge L_0$ and every parameter $s$.

An infinite-volume GNS witness instead chooses a ground state $\omega_s$,
its GNS representation, and a positive generator $K_s$ of the implemented
dynamics such that
\begin{equation}
\Spec(K_s)\cap(0,\gamma)=\varnothing
\label{eq:gns-gap}
\end{equation}
uniformly in $s$.  Relating the two notions requires hypotheses on boundary
conditions, ground-state sectors, and the thermodynamic limit.  The paper
does not treat them as definitionally equal.

\subsection{Witnesses form a fibration of data}

Let $W(S)$ denote the type of accepted gap witnesses over an $S$-family.
Define
\begin{equation}
\Gap^{\mathrm{wit}}_{\F}(S)
=
\{(H,w): H\in\Ham_{\F}(S),\ w\in W_S(H)\}.
\label{eq:witness-gap-object}
\end{equation}
Forgetting the numerical and ground-state data gives
\begin{equation}
p_S:\Gap^{\mathrm{wit}}_{\F}(S)\longrightarrow\Ham_{\F}(S).
\label{eq:witness-projection}
\end{equation}

\begin{proposition}[Witness projection is not a subobject]
\label{prop:witness-not-subobject}
Even when all restriction maps are defined, $p_S$ need not be a
monomorphism.  One Hamiltonian family can admit many gap lower bounds,
thresholds, ground-state presentations, and equivalent proofs.  Thus
$\Gap^{\mathrm{wit}}$ is a fibration of quantitative data over the
Hamiltonian moduli object, not the qualitative gapped substack.
\end{proposition}

This is more than a set-theoretic nuisance.  Loops in the space of witnesses
can contain information about choices that should disappear from a phase
label.  Treating \cref{eq:witness-projection} as an inclusion would retain
that information by mistake.

\subsection{The propositional image}

The qualitative predicate is obtained by truncating the existence of a
witness:
\begin{equation}
\Gap^{\mathrm{prop}}_{\F}(S)
=
\{H\in\Ham_{\F}(S):\|\exists w\in W_S(H)\|\}.
\label{eq:propositional-gap}
\end{equation}
Here the vertical bars denote propositional truncation.  Equivalently,
$\Gap^{\mathrm{prop}}$ is the image of $p$ when the chosen category of
condensed anima has pullback-stable image factorizations.

\begin{definition}[Uniformly gapped propositional substack, proposed]
\PROP{}  If gap witnesses restrict functorially, compatible witnesses glue
over the chosen covers with a common positive lower bound, and the image
factorization in \cref{eq:propositional-gap} commutes with pullback, then
$\Gap^{\mathrm{prop}}_{\F}\hookrightarrow\Ham_{\F}$ is the uniformly gapped
propositional substack.
\end{definition}

Descent can combine witnesses that already exist.  It cannot infer a common
positive lower bound from unrelated local or pointwise gaps.  For a finite
cover, finitely many compatible lower bounds have a positive minimum.  For
an infinite cover, the infimum can be zero.

\subsection{Gap stability}

\begin{theorem}[Controlled stability, schematic]
\label{thm:controlled-gap-stability}
\EST{} in its cited operator-algebraic settings.  Let
$r\mapsto\Phi_r$ be a path of interactions with common locality bounds and a
uniform thermodynamic gap.  Under the standard differentiability and decay
hypotheses, quasi-adiabatic spectral flow gives quasi-local automorphisms
that transport the ground-state sectors along the path.  Phase invariants
defined through this transport remain constant
\cite{BMNS2012,BachmannNachtergaele2014}.
\end{theorem}

The theorem begins with a uniform gap.  It does not prove that a proposed
path is gapped.  Establishing the gap is often the hardest analytic task in
a phase classification.

\section{Stage four: the stabilized phase infinity-groupoid}

The qualitative gapped locus still contains every microscopic
presentation.  A phase object should identify systems connected by the
allowed physical moves.

\subsection{The path category}

For a test object $S$, let $\C_{\mathrm{gap}}(S)$ be the proposed category
whose objects are $S$-families in $\Gap^{\mathrm{prop}}(S)$ and whose
morphisms are controlled uniformly gapped interpolations, quasi-local
automorphisms supplied by spectral flow, and specified changes of local
presentation.  Every morphism carries a provider showing that locality and
the gap remain uniform over both $S$ and the interpolation parameter.

Let $W(S)$ be the class of morphisms declared to be phase equivalences.  A
calculus of fractions is not assumed automatically.  The infinity-categorical
localization
\begin{equation}
\C_{\mathrm{gap}}(S)[W(S)^{-1}]
\label{eq:phase-localization}
\end{equation}
exists abstractly in a sufficiently large ambient category, but size,
descent, and compatibility with stacking must still be checked.

\subsection{Stabilization}

Choose a class $\mathsf{At}$ of atomic product-state systems.  Stacking with
an atomic ancilla gives transition functors
\begin{equation}
H\longmapsto H\boxtimes A,
\qquad A\in\mathsf{At}.
\end{equation}
The stabilized category is a filtered colimit only after maps between ancilla
choices and their coherences have been specified.  Formally, write
\begin{equation}
\C_{\mathrm{gap}}^{\mathrm{st}}(S)
=
\operatorname*{colim}_{A\in\mathsf{At}}
\C_{\mathrm{gap}}(S)[W(S)^{-1}].
\label{eq:stabilized-category}
\end{equation}

\begin{definition}[Stabilized phase infinity-groupoid, proposed]
\PROP{}  Under the existence and descent hypotheses above, set
\begin{equation}
\Phase^{\mathrm{st}}_{d,\F,G}(S)
=
\Core\bigl(\C_{\mathrm{gap}}^{\mathrm{st}}(S)\bigr).
\label{eq:phase-groupoid}
\end{equation}
The result is intended to be an anima-valued condensed sheaf.
\end{definition}

Taking the core discards noninvertible categorical morphisms but retains all
equivalences and their higher homotopies.  It does not make every object
invertible under stacking.  That is a separate operation.

\subsection{Why stabilization changes the question}

Two finite-band Hamiltonians may be inequivalent before adding trivial
bands and equivalent afterward.  Two spin chains may become equivalent
after stacking with a product state that changes the on-site Hilbert space.
Stabilization is therefore a quotient choice, not a harmless notational
convenience.

The BDI example goes further.  Allowing symmetry-preserving interactions
enlarges the path category and reduces the free integer classification to
$\mathbb Z/8$ \cite{FidkowskiKitaev2010}.  A stable free K-class does not survive unchanged when the
microscopic morphisms are changed.

\section{Stage five: the invertible condensed phase spectrum}

Stacking supplies a symmetric monoidal product on microscopic systems.  If
it descends through the gap condition, localization, and stabilization, it
makes $\Phase^{\mathrm{st}}$ an $E_\infty$ object in condensed anima.

\subsection{The invertible sector}

An object $x$ is stacking-invertible when there is an object $y$ and an
equivalence
\begin{equation}
x\boxtimes y\simeq\mathbf 1
\end{equation}
in the stabilized phase infinity-groupoid.  The full subobject of such
objects is
\begin{equation}
\Pic(\Phase^{\mathrm{st}}).
\label{eq:picard-phase}
\end{equation}
It is grouplike if the monoidal coherences and descent are in place.

\begin{remark}[Group completion and units]
One may group-complete the entire stacking monoid and then study its units,
or take the already invertible full subobject.  These operations need not
encode the same physical question if formal inverses exist without
microscopic witnesses.  The program here uses the Picard subobject: an
inverse must be represented by a phase.
\end{remark}

\subsection{Spectrification}

Grouplike $E_\infty$ spaces are equivalent to connective spectra.  In a
suitable infinity-topos, the same recognition principle can be applied
internally, subject to presentability and descent hypotheses.  This motivates
the definition
\begin{equation}
\IP^{\mathrm{cond}}_{d,\F,G}
=
\operatorname{sp}
\Pic(\Phase^{\mathrm{st}}_{d,\F,G}).
\label{eq:condensed-spectrum}
\end{equation}

\begin{definition}[Invertible condensed phase spectrum, proposed]
\PROP{}  Equation \cref{eq:condensed-spectrum} denotes the connective
condensed spectrum obtained from the stacking-invertible stabilized phase
object, provided the latter exists as a grouplike $E_\infty$ condensed anima
and the internal recognition theorem applies in the chosen category.
\end{definition}

This definition does not prove that the spectrum exists for the full class
of interacting lattice systems.  It also does not identify its spatial
degree $d$ pieces with the levels of an Omega-spectrum.  A spatial suspension
or lattice-direction delooping theorem is still required.

\subsection{A conditional assembly theorem}

The logical content of the program can be packaged as follows.

\begin{theorem}[Conditional assembly]
\label{thm:conditional-assembly}
Assume:
\begin{enumerate}[label=(A\arabic*),leftmargin=2.7em]
  \item the uniformly $F$-local interaction presheaf has functorial
  restrictions and hyperdescent;
  \item Hamiltonian objects, controlled morphisms, and higher equivalences
  form an anima-valued condensed stack;
  \item quantitative finite-volume or GNS gap witnesses restrict and descend
  with a common positive lower bound;
  \item the witness projection admits a pullback-stable propositional image;
  \item controlled gapped paths and quasi-local equivalences admit a
  localization compatible with descent;
  \item stabilization by the declared atomic systems exists and commutes with
  restriction;
  \item stacking is coherently symmetric monoidal throughout the
  construction and preserves all declared equivalences;
  \item the Picard subobject is a grouplike $E_\infty$ condensed anima to
  which internal connective spectrification applies.
\end{enumerate}
Then the five stages define a connective invertible condensed phase spectrum
$\IP^{\mathrm{cond}}_{d,\F,G}$.  Its zeroth homotopy sheaf is the sheaf of
stacking-invertible stabilized phases represented in the construction.
\end{theorem}

\begin{proof}
Assumptions (A1) and (A2) give the condensed Hamiltonian moduli object.
Assumptions (A3) and (A4) give a propositional gapped subobject without
retaining choices of gap proof.  Assumptions (A5) and (A6) produce the
stabilized phase infinity-groupoid.  Assumption (A7) equips it with a
coherent $E_\infty$ stacking product.  Assumption (A8) identifies the Picard
subobject with the zeroth space of a connective spectrum internal to the
chosen condensed setting.  The description of the zeroth homotopy sheaf is
the definition of connected components in that Picard object.
\end{proof}

The theorem is formally valid but conditional.  Its conclusion should not
be cited without the assumptions.  Current work establishes parts of (A1),
(A3), and (A5) in controlled settings.  The full stack-level assembly is
\PROP{}, and broad comparison or spatial delooping statements are \OPEN{} or
\CONJ{}.

\section{Homotopy groups, pumps, and defects}

The slogan
\begin{align}
\pi_0&=\text{phases},\notag\\
\pi_1&=\text{adiabatic pumps and phase automorphisms},\notag\\
\pi_k&=\text{higher families and higher defects}
\label{eq:homotopy-slogan}
\end{align}
is useful only after the basepoint and the kind of homotopy group have been
specified.

\subsection{Homotopy sheaves and global shapes}

For a based condensed anima $X$, one can form homotopy sheaves
$\underline\pi_k(X)$ on test objects.  One can also take global sections or
a shape and then form ordinary homotopy groups.  These operations need not
agree.  In the proposed spectrum, the least ambiguous object is the homotopy
sheaf
\begin{equation}
\underline\pi_k\bigl(\IP^{\mathrm{cond}}_{d,\F,G}\bigr).
\end{equation}
Its value at a test object $S$ records $S$-parameterized classes, subject to
the success of the assembly theorem.

\subsection{Zeroth homotopy}

At a point test object,
\begin{equation}
\pi_0\Omega^\infty\IP^{\mathrm{cond}}
\end{equation}
would be the abelian group of stacking-invertible stabilized phases.  It is
not the set of all topological phases, since noninvertible phases have been
removed.  It also depends on symmetry, allowed interactions, boundary
conventions, and the selected stabilization.

\subsection{First homotopy}

A based loop in the phase infinity-groupoid is a one-parameter family that
returns to its starting phase up to the declared equivalence.  With a
uniform gap and a quasi-adiabatic transport provider, such a loop can define
an adiabatic pump or an automorphism of the phase.  Without that analytic
provider it is only a loop in the proposed moduli object.

Thus the identification of $\pi_1$ with pumps has two layers:
\begin{enumerate}[leftmargin=2.2em]
  \item categorical loop data in the stabilized phase object;
  \item an analytic transport theorem that assigns a physical pumped
  quantity or boundary action.
\end{enumerate}
The first is part of the proposed organization.  The second is established
in important model classes and remains open in full generality.

\subsection{Higher families}

A class in $\pi_k$ can be represented by a based $k$-parameter family.  Its
interpretation as a codimension-$(k+1)$ defect requires a separate
family-to-defect construction, usually involving suspension, boundary
conditions, or a clutching map.  Higher homotopy does not turn into a defect
merely by changing terminology.

\begin{proposition}[Qualified higher-defect interpretation]
Suppose a family-to-defect map is defined on controlled uniformly gapped
families, respects homotopy and stacking, and is compatible with the chosen
spatial suspension.  Then a class in the relevant $\pi_k$ determines a
stable defect class.  Without this map, $\pi_k$ records a higher family but
not yet a physical defect.
\end{proposition}

\section{The gapless locus and relative transition charge}

Let $B$ be a parameter space and let
\begin{equation}
f:B\longrightarrow\Ham_{\F}
\label{eq:parameter-map}
\end{equation}
be a Hamiltonian family.  Pull back the propositional gapped subobject:
\begin{equation}
U_f
=
B\times_{\Ham_{\F}}\Gap^{\mathrm{prop}}_{\F}.
\label{eq:gapped-pullback}
\end{equation}
When $U_f$ is represented by an open subspace of $B$, define
\begin{equation}
\Sigma_f=B\setminus U_f.
\label{eq:gapless-locus}
\end{equation}
The symbol $\Sigma$ therefore denotes a qualitative failure locus.  The
witness fibration also pulls back to $U_f$, but it cannot be subtracted from
$B$ because it is not a subobject.  If openness, closedness of $\Sigma$, or a
good-pair replacement is unavailable, the relative group used below is not
assigned automatically to the proposed gapless locus.

\subsection{Why openness is an analytic theorem}

For bounded self-adjoint elements, a spectral gap is stable under sufficiently
small norm perturbations.  For thermodynamic interacting systems, openness
requires a topology on interactions and a stability theorem uniform in
volume.  The statement ``the gapped locus is open'' is therefore established
only in a declared controlled setting.  Condensed sheafification alone does
not prove it.

\subsection{The exact relative class}

Let $E$ be a generalized cohomology theory represented by a spectrum.  Put
$U=B\setminus\Sigma$ and assume $(B,U)$ is a good pair for $E$.  The long
exact sequence contains
\begin{equation}
E^q(B)\xrightarrow{j^*}E^q(U)
\xrightarrow{\partial}E^{q+1}(B,U)
\longrightarrow E^{q+1}(B).
\label{eq:relative-les}
\end{equation}
For $\nu\in E^q(U)$, define
\begin{equation}
Q_\Sigma(\nu)=\partial\nu
\in E^{q+1}(B,B\setminus\Sigma).
\label{eq:relative-charge}
\end{equation}

\begin{proposition}[Extension obstruction]
\label{prop:extension-obstruction}
\EST{}  The class $Q_\Sigma(\nu)$ vanishes exactly when $\nu$ lies in the
image of $E^q(B)\to E^q(U)$.
\end{proposition}

\begin{proof}
This is exactness of \cref{eq:relative-les} at $E^q(U)$.
\end{proof}

The proposition concerns extension of a cohomology class.  It does not say
that the Hamiltonian family extends across $\Sigma$ as a uniformly gapped
family.

\subsection{Excision and Thom localization}

If excision applies near the closed set $\Sigma$, then
\begin{equation}
E^{q+1}(B,B\setminus\Sigma)
\cong
E^{q+1}(N,N\setminus\Sigma)
\label{eq:excision-local}
\end{equation}
for a suitable neighborhood $N$.  Suppose further that $B$ is smooth,
$\Sigma\hookrightarrow B$ is a closed submanifold of codimension $c$, and
its normal bundle $\mathcal N$ is $E$-oriented.  A tubular neighborhood and
the Thom isomorphism give
\begin{equation}
E^{q+1}(B,B\setminus\Sigma)
\cong
\widetilde E^{q+1}(\Th(\mathcal N))
\cong
E^{q+1-c}(\Sigma).
\label{eq:thom-charge}
\end{equation}
Without an $E$-orientation, the last group must be replaced by the
appropriate twisted theory.  For singular $\Sigma$, a stratified normal
datum or supported-cohomology formulation is needed.

\subsection{When the class is physical}

Suppose a microscopic family $H$ on $U$ has a phase invariant
$\nu_{\mathrm{mic}}(H)$.  To obtain the $E$-cohomology class used above one
needs a natural transformation
\begin{equation}
c_E:\nu_{\mathrm{mic}}(H)\longmapsto\nu_E(H)\in E^q(U)
\label{eq:physical-comparison}
\end{equation}
that respects restriction, stacking, and controlled gapped homotopy.

\begin{theorem}[Conditional physical transition charge]
\label{thm:conditional-charge}
If \cref{eq:physical-comparison} is supplied and has the stated naturality,
then
\begin{equation}
Q_\Sigma(H):=\partial\nu_E(H)
\end{equation}
is invariant under the chosen microscopic phase relation on $U$.  It is
localized near $\Sigma$ by excision and, under the Thom hypotheses, has
degree $q+1-c$ on the critical locus.
\end{theorem}

The algebraic-topology part of this theorem is \EST{}.  The general
microscopic comparison is \PROP{} or \OPEN{}, depending on the model class.
For a controlled free-fermion family, functional calculus can provide it.
For arbitrary interacting systems, no universal map is asserted.

\subsection{A local SSH test}

For the chiral two-band SSH Hamiltonian, the off-diagonal function
\begin{equation}
q(k)=t_1+t_2e^{ik}
\end{equation}
has winding one when $|t_1|<|t_2|$ and winding zero when
$|t_1|>|t_2|$.  At the isolated crossing $t_1=t_2$, $k=\pi$, a small linking
circle in the combined momentum and mass coordinates has degree of magnitude
one.  This is a controlled example of \cref{eq:relative-charge}.  A loop in
the tuning parameter alone would miss the momentum coordinate and need not
link the degeneracy \cite{TeoKane2010,Thiang2016}.

\section{Microscopic systems and effective field theories}

The phrase ``microscopic to effective'' covers three separate operations.
They should be drawn as three arrows:
\begin{equation}
\begin{tikzcd}[column sep=huge,row sep=large]
\Ham^{\mathrm{micro,gap}}
  \arrow[r,dashed,"L"]
& \mathsf{EFT}^{\mathrm{inv}}
  \arrow[r,"I","\text{established on a stated EFT domain}"']
& \mathsf{Class}^{\mathrm{stable}}
  \arrow[dll,dashed,"R"] \\
\mathsf{Real}_{\mathrm{lattice}}
  \arrow[u,"\text{forget witnesses}"']
&&
\end{tikzcd}
\label{eq:three-comparisons}
\end{equation}
Here $L$ is a low-energy limit, $I$ extracts a stable invariant on a specified
effective-theory domain, and $R$ is a
realization problem whose target contains explicit stabilized microscopic
witnesses.  The final vertical map forgets those witnesses and retains the
underlying microscopic Hamiltonian.  None is automatically inverse to
another.  The solid $I$ does not assert an invariant map on every object
called an EFT; its established use here is restricted to declared targets,
including the Freed-Hopkins domain described below.

\subsection{Controlled free functional calculus}

Let $h\in M_N(C(S,\A))_{\mathrm{sa}}$ and suppose
\begin{equation}
\Spec(h_s)\cap(-\gamma,\gamma)=\varnothing
\end{equation}
for a common $\gamma>0$.  Continuous functional calculus gives
\begin{equation}
q_h=\operatorname{sgn}(h),
\qquad
p_h=\frac{1-q_h}{2}.
\label{eq:flattening}
\end{equation}
The projection $p_h$ varies continuously with $s$.  Relative to a reference
$h_{\mathrm{ref}}$, the difference
\begin{equation}
\kappa(h,h_{\mathrm{ref}})
=[p_h]-[p_{h_{\mathrm{ref}}}]
\label{eq:k-class}
\end{equation}
is invariant under norm-continuous paths with a common gap and is additive
under block sum.  Throughout this paper the first argument is the target and
the second is the reference, so $\kappa(\text{target},\text{reference})$ means
target minus reference.

This is an \EST{} map in the controlled free sector.  It loses the gap
magnitude, dispersion, Fermi velocity, correlation length, Lieb-Robinson
constants, coupling values, boundary termination, and much of the band
presentation.  Flattening is not a renormalization-group construction.

\subsection{Interacting systems}

A general interacting many-body Hamiltonian does not have a bounded
one-particle matrix to which \cref{eq:flattening} applies.  Flattening a
finite-volume matrix can produce a highly nonlocal operator and may not have
a quasi-local thermodynamic limit.  Quasi-adiabatic spectral flow is the
right comparison tool for paths of interacting systems, but it starts from
uniform locality and a uniform gap.

The Fidkowski-Kitaev reduction is a decisive warning.  One-dimensional BDI
free phases have an integer index.  Symmetry-preserving interactions reduce
the classification to $\mathbb Z/8$ \cite{FidkowskiKitaev2010}.  Therefore no universal equivalence can
identify free operator K-theory with the interacting microscopic phase
object.

\subsection{Disorder}

For a covariant disordered free Hamiltonian with a genuine C*-spectral gap,
the Fermi projection lies in an appropriate crossed-product algebra and
defines a K-class.  A mobility-gap extension needs localization estimates
strong enough to place the projection and its commutators in the relevant
algebra.  A profinite disorder hull does not supply those estimates.

\section{Three stable-homotopy targets that remain distinct}

The proposed condensed spectrum is not the only spectrum in the subject.
Two established constructions have different domains.

\subsection{Freed-Hopkins field-theory classes}

Fix a spacetime symmetry type $(H_n,\rho_n)$ and its Madsen-Tillmann spectrum
$MTH$.  In the discrete topological sector, Freed and Hopkins compute
deformation classes of reflection-positive invertible extended topological
field theories by a torsion subgroup of
\begin{equation}
[MTH,\Sigma^{n+1}I\mathbb Z(1)].
\label{eq:fh-target}
\end{equation}
This is an \EST{} field-theory theorem under its hypotheses
\cite{FreedHopkins2021}.  Application to
lattice phases uses an effective-field-theory assumption.  The theorem does
not assert that every class in \cref{eq:fh-target} has a finite-range lattice
representative with a thermodynamic gap.

\subsection{Kubota's microscopic Omega-spectrum}

Kubota constructs an \EST{} Omega-spectrum $\IP_*$ from invertible gapped
quantum spin systems on Euclidean spaces.  Its ingredients include
operator-algebraic interactions, uniformly almost-local control, relaxed
lattices, stabilization by atomic systems, nondegenerate gapped GNS ground
states, and sheaves of smooth parameter families \cite{Kubota2025}.  This is a genuine
microscopic stable-homotopy object.

It is not definitionally the condensed spectrum in
\cref{eq:condensed-spectrum}.  The parameter sites, gap witnesses, lattice
conventions, and descent conditions differ.  A comparison should first be
defined on their common smooth, uniformly controlled domain.

\subsection{Aoki's K-theory bridge}

Aoki proves, with connective and Bott-periodic qualifications, that
solidification of algebraic K-theory of real or complex Banach algebras
recovers the corresponding connective topological K-theory, followed by
inversion of the degree-two complex or degree-eight real Bott element for
the periodic theory.  This is an \EST{} invariant-level bridge
\cite{Aoki2024}.

The input Banach or C*-algebra is already present.  Solidification does not
recover its norm, positive cone, state space, or microscopic interaction.
It also does not say that every K-class is a physical phase.

\subsection{Comparison diagram}

The honest current diagram is
\begin{equation}
\begin{tikzcd}[column sep=large,row sep=large]
& \IP_*^{\mathrm{Kubota}}
  \arrow[dl,dotted,"C_{\mathrm{cond}}"']
  \arrow[dr,dotted,"C_{\mathrm{EFT}}"] & \\
\IP^{\mathrm{cond}}
  \arrow[rr,dotted,"C_{\mathrm{bord}}"']
&& \Map(MTH,\Sigma^{*+1}I\mathbb Z(1)) \\
& \Kop
  \arrow[ul,dotted,"C_K"]
  \arrow[ur,dotted,"C_{K,\mathrm{EFT}}"'] &
\end{tikzcd}
\label{eq:comparison-web}
\end{equation}
Every dotted arrow needs a domain, a formula, naturality, stacking
compatibility, and homotopy invariance.  An equivalence claim also needs
full faithfulness and essential surjectivity, or their spectral analogues.
No such global equivalence is asserted here.

\section{Physical realizability}

An abstract class becomes physical only through a witness.  For a
microscopic system in spatial dimension $d$, a realization record contains:
\begin{enumerate}[leftmargin=2.2em]
  \item local degrees of freedom and a lattice;
  \item an explicit uniformly local interaction;
  \item symmetry, grading, and boundary conventions;
  \item a thermodynamic gap witness;
  \item a microscopic invariant;
  \item an explicit comparison to the abstract class;
  \item a stacking inverse if invertibility is claimed.
\end{enumerate}
An element of a bordism group supplies none of the microscopic fields by
itself.

\subsection{Three realized low-dimensional rows}

The companion realizability analysis establishes three concrete rows.
\begin{longtable}{>{\raggedright\arraybackslash}p{0.16\textwidth}
                  >{\raggedright\arraybackslash}p{0.20\textwidth}
                  >{\raggedright\arraybackslash}p{0.23\textwidth}
                  >{\raggedright\arraybackslash}p{0.30\textwidth}}
\toprule
System & Gap witness & Microscopic invariant & Abstract comparison status \\
\midrule
\endfirsthead
\toprule
System & Gap witness & Microscopic invariant & Abstract comparison status \\
\midrule
\endhead
Class-D Kitaev chain
& Exact flat periodic bulk spectrum at the solvable point
& Pfaffian parity and one Majorana at each open end
& Agreement with the nontrivial spin or Arf $\mathbb Z/2$ class under the low-energy interpretation \cite{Kitaev2001,FreedHopkins2021} \\

BDI stack
& Exact free gap plus the Fidkowski-Kitaev gapped interaction for eight copies
& Free integer, reduced to residue modulo eight with interactions
& Agreement with the two-dimensional $\mathrm{Pin}^-$ Arf-Brown-Kervaire $\mathbb Z/8$ sector \cite{FidkowskiKitaev2010,FreedHopkins2021} \\

Spin-one AKLT chain
& Rigorous AKLT bulk-gap theorem
& Nontrivial projective $SO(3)$ boundary representation
& Agreement with the nontrivial $H^2(SO(3),U(1))$ sector \cite{AKLT1987,BachmannNachtergaele2014} \\
\bottomrule
\end{longtable}

Each row is microscopically realized in the stated one-dimensional setting.
The last column records a limited comparison, not a spectrum equivalence.

\subsection{Surjectivity and injectivity}

Surjectivity of a realization map would require every abstract class to have
a local Hamiltonian witness with a uniform gap and the correct response.
Anomalies, finite-dimensional local Hilbert spaces, spatial symmetry, or
unknown gap proofs may obstruct or delay such a result.

Injectivity would require two systems with the same abstract class to be
connected by the chosen stabilized gapped path.  A field-theory target may
forget crystalline data, noninvertible excitations, boundary termination,
or unstable finite-band structure.  If the microscopic category retains any
of these, injectivity can fail.

The BDI example shows that changing the domain can also merge classes.  Once
interactions are admitted, paths exist that were absent in the free
category.  Realization and comparison questions are meaningless until the
domain and its morphisms are fixed.

\section{Arrow-by-arrow theorem obligations}

The five-stage program is best audited one arrow at a time.

\begin{longtable}{>{\raggedright\arraybackslash}p{0.19\textwidth}
                  >{\raggedright\arraybackslash}p{0.31\textwidth}
                  >{\raggedright\arraybackslash}p{0.15\textwidth}
                  >{\raggedright\arraybackslash}p{0.24\textwidth}}
\toprule
Arrow & Required theorem & Status & Main missing point \\
\midrule
\endfirsthead
\toprule
Arrow & Required theorem & Status & Main missing point \\
\midrule
\endhead
Interactions to Hamiltonian moduli
& Uniform analytic bounds define a stack with controlled morphisms
& \PROP{}
& Higher descent and variable microscopic presentations \\

Hamiltonian moduli to gapped locus
& Witness projection has pullback-stable propositional image
& \PROP{}
& Uniform thermodynamic witnesses and image descent \\

Gapped locus to stabilized phases
& Controlled path localization and atomic stabilization commute with descent
& \PROP{}
& Global calculus of equivalences and coherences \\

Stabilized phases to invertible sector
& Stacking descends as a symmetric monoidal structure and units form a Picard object
& \PROP{}
& Physical inverses and sheafwise invertibility \\

Picard object to spectrum
& Internal grouplike recognition and spatial delooping
& \PROP{} for connective spectrification, \OPEN{} for spatial Omega-structure
& Compatibility between spatial suspension and categorical delooping \\

Microscopic spectrum to K-theory
& Natural invariant map preserving families and stacking
& \EST{} in controlled free sectors, \OPEN{} generally
& Interactions and information loss \\

Microscopic spectrum to EFT or bordism
& Low-energy functor with stabilization and defect compatibility
& \OPEN{} generally
& Existence, functoriality, faithfulness, and realization \\

Kubota spectrum to condensed spectrum
& Comparison on a common smooth and condensed parameter domain
& \OPEN{}
& Site change, gap conventions, and descent \\
\bottomrule
\end{longtable}

This table prevents progress on one arrow from being reported as completion
of the sequence.  Aoki's theorem, for example, advances an invariant bridge
after a Banach algebra is given.  It does not solve the moduli, gap, phase,
or realization arrows.

\section{Obstructions and failed shortcuts}

Several tempting shortcuts are known to fail.

\begin{counterexample}[Pointwise gaps]
For every parameter $s$ and every finite volume, suppose the Hamiltonian has
a positive gap $\gamma_{s,L}$.  The infimum over $s$ and $L$ may still be
zero.  The data do not define a point of $\Gap^{\mathrm{wit}}$ without one
positive lower bound and a declared low-energy sector.
\end{counterexample}

\begin{counterexample}[Qualitative image replaced by witnesses]
If a Hamiltonian admits two lower bounds $\gamma$ and $\gamma/2$, the witness
object contains at least two points over the same Hamiltonian.  Its projection
cannot be the inclusion of a substack.  The qualitative image must truncate
the choice.
\end{counterexample}

\begin{counterexample}[Homotopy groups without a map]
Two connective spectra can have isomorphic individual homotopy groups while
their Postnikov data or multiplicative structures differ.  A list of group
agreements does not define an equivalence of spectra.
\end{counterexample}

\begin{counterexample}[Finite-size numerical gap evidence]
A computation at lengths $L\le20$ can be consistent with a sequence of gaps
that later tends to zero.  It is evidence for a conjecture, not a
thermodynamic witness.
\end{counterexample}

\begin{counterexample}[Free invariant under interacting paths]
BDI phases of free indices zero and eight are distinct in the free path
category and equivalent after the Fidkowski-Kitaev interaction is admitted.
Thus the free K-class is not a complete invariant of the interacting phase
category.
\end{counterexample}

\begin{counterexample}[Bordism generator without lattice data]
A generator of a bordism or Anderson-dual group specifies an abstract
field-theory deformation class.  It does not specify local degrees of
freedom, an interaction, a gap theorem, or a stacking inverse.  Calling the
class realized before those witnesses are supplied changes the meaning of
realization.
\end{counterexample}

\section{A precise research agenda}

The synthesis isolates five theorem packages.  They can be pursued largely
independently and then joined by explicit interfaces.

\subsection{Package I: analytic family control}

The first package should formulate a category of lattices and interactions
on which a common locality norm implies family-uniform Lieb-Robinson bounds,
thermodynamic dynamics, and restriction stability.  It should include
moving defects, disorder hulls, and variable local Hilbert spaces without
confusing local continuity with uniform decay.

The main target is a base-change theorem: pullback of a controlled family
preserves every constant used by the dynamics.  A second target is finite
descent for quantitative locality certificates.  Infinite descent will
need a bornological or boundedness condition that prevents constants from
escaping.

\subsection{Package II: gap certificates and openness}

The second package should compare finite-volume and GNS witnesses under
explicit boundary and sector hypotheses.  Its output should be a stack of
witnesses over the Hamiltonian moduli object, together with a theorem that
the propositional image is pullback-stable.

A useful certificate has fields for the lower bound, volume threshold,
ground-state sector, boundary convention, and theorem provider.  This makes
the gap claim transportable without pretending that the provider is unique.

\subsection{Package III: phase localization}

The third package should define controlled gapped paths, quasi-local
automorphisms, finite-depth circuits, lattice changes, and stabilization in
one infinity-categorical model.  Comparison functors between these notions
should be theorems, not definitional equalities.

The practical question is whether local phase equivalences glue.  On an
overlap, two spectral-flow providers may differ by a phase automorphism.
That difference belongs in the higher groupoid and cannot be erased during
descent.

\subsection{Package IV: stacking and spectrification}

The fourth package should prove that stacking preserves locality and gap
witnesses with controlled constants, descends through localization, and is
coherently symmetric monoidal.  It should then distinguish actual physical
inverses from formal Grothendieck inverses.

Connective spectrification is only part of the job.  To compare spatial
dimensions, one needs a theorem that adding a lattice direction implements
the loop or suspension maps of an Omega-spectrum.  Kubota's construction is
the clearest established benchmark for this step.

\subsection{Package V: comparison and realization}

The fifth package should build natural transformations on carefully chosen
overlaps:
\begin{enumerate}[leftmargin=2.2em]
  \item from free condensed Hamiltonian families to operator K-theory;
  \item from Kubota's smooth microscopic families to condensed tests;
  \item from controlled microscopic invertible phases to reflection-positive
  effective field theories;
  \item from a generalized cohomology invariant to relative transition
  classes;
  \item from abstract classes back to explicit microscopic witnesses where
  possible.
\end{enumerate}
Each map should list the information it forgets.  A realization theorem
should state its locality, symmetry, dimensional, and boundary restrictions.

\section{Formal and executable representations}

The repository contains a Lean interface library for the shared vocabulary
of the project.  It records claim statuses, providers, locality and gap
witness interfaces, objectwise C*-conditions, stacking, phase equivalence,
and comparison records.  It does not assert the deep analytic theorems as
axioms.  This synthesis does not modify that library.

The Haskell program accompanying this paper has a narrower purpose.  It
checks the architecture of a proposed synthesis record.  Its data types
distinguish the witness fibration from the propositional image and distinguish
an invariant map from an equivalence theorem.  Its tests reject:
\begin{enumerate}[leftmargin=2.2em]
  \item a phase sequence that skips a stage or reverses its order;
  \item a qualitative gapped locus represented by witness data;
  \item an equivalence claim without a named comparison theorem;
  \item a physical relative charge without a microscopic comparison;
  \item an untwisted Thom target without an orientation;
  \item a spacetime degree that fails $n=d+1$.
\end{enumerate}

The executable is not a proof assistant.  It verifies a finite status and
dependency contract.  Analytic results remain cited theorem providers.

\section{Claim ledger}

\begin{longtable}{>{\raggedright\arraybackslash}p{0.24\textwidth}
                  >{\raggedright\arraybackslash}p{0.15\textwidth}
                  >{\raggedright\arraybackslash}p{0.52\textwidth}}
\toprule
Statement & Status & Basis and boundary \\
\midrule
\endfirsthead
\toprule
Statement & Status & Basis and boundary \\
\midrule
\endhead
Uniform Lieb-Robinson bounds and thermodynamic dynamics
& \EST{}
& Common interaction norm, bounded geometry, and uniform tail hypotheses \\

C*-norm and positivity for $C(S,\A)$
& \EST{}
& A fixed C*-algebra $\A$ and compact Hausdorff parameter space \\

Finite-cover C*-descent
& \EST{}
& Finite jointly surjective profinite cover and continuous gluing \\

Bare algebraic data determine a unique physical C*-norm
& \OBST{}
& Distinct C*-completions provide counterexamples \\

Uniform gap stability and spectral flow
& \EST{}
& Controlled paths satisfying the established locality, regularity, and uniform-gap hypotheses \\

Pointwise gaps imply a family-uniform thermodynamic gap
& \OBST{}
& Positive numbers can have zero infimum; moving defects give physical models \\

Witnessed gap object projects as a fibration
& \PROP{}
& Explicit dependent-pair construction with many witnesses per Hamiltonian \\

Propositional gapped substack
& \PROP{}
& Requires pullback-stable image factorization and descent of common witnesses \\

Stabilized phase infinity-groupoid
& \PROP{}
& Requires controlled localization, ancilla stabilization, and higher descent \\

Invertible condensed phase spectrum
& \PROP{}
& Conditional internal spectrification of the Picard phase object \\

Spatial Omega-spectrum structure for the condensed construction
& \OPEN{}
& Needs a suspension or lattice-direction delooping theorem \\

Controlled free flattening to operator K-theory
& \EST{}
& Bounded self-adjoint family with a common Fermi gap \\

Aoki solidification comparison
& \EST{}
& Algebraic to topological K-theory for real or complex Banach algebras, with connective and Bott qualifications \\

Free K-theory classifies all interacting phases
& \OBST{}
& BDI interaction reduction supplies a counterexample \\

Kubota microscopic Omega-spectrum
& \EST{}
& Invertible gapped spin systems with the hypotheses of Kubota's construction \\

Freed-Hopkins field-theory classification
& \EST{}
& Reflection-positive invertible field theories under the stated tangential and field-theoretic hypotheses \\

Kubota spectrum equals the condensed phase spectrum
& \OPEN{}
& No comparison equivalence has been constructed \\

Microscopic spectrum equals the Freed-Hopkins target
& \OPEN{}
& Low-energy functor, family compatibility, and realization are missing generally \\

Kitaev, BDI, and AKLT witness rows
& \EST{}
& Explicit interactions, symmetries, gap calculations or theorems, and microscopic invariants \\

Every bordism class is physically realizable
& \OPEN{}
& No general surjectivity theorem is asserted \\

Relative connecting class $\partial\nu$
& \EST{}
& Long exact sequence of the pair $(B,B\setminus\Sigma)$ \\

Thom localization to $E^{q+1-c}(\Sigma)$
& \EST{}
& Excision, smooth closed locus, and $E$-oriented rank-$c$ normal bundle \\

Every relative class is a physical transition charge
& \OBST{}
& A microscopic invariant and natural comparison map are additional required data \\
\bottomrule
\end{longtable}

\section{Discussion}

The condensed perspective earns its keep by forcing parameter dependence to
be part of the object.  Disorder, inverse limits, and higher families then
belong to the same moduli problem as ordinary paths.  It also forces descent
questions into the open.  A family assembled from local pieces needs common
analytic constants, not just compatible point values.

The separation of witnessed and qualitative gaps is indispensable.  A gap
proof contains a lower bound, a ground-state convention, and often a
theorem-specific presentation.  Those choices matter during verification
and are irrelevant to the yes-or-no gapped locus.  The witness fibration
retains them; the propositional image forgets them.  Conflating the two
damages both the mathematics and the software model.

Stable-homotopy objects in this subject also come from different physical
categories.  Operator K-theory is an invariant of a specified
operator algebra.  Kubota's spectrum is built from microscopic gapped spin
systems and smooth parameter families.  Freed-Hopkins classify invertible
field theories under reflection positivity and tangential hypotheses.  The
proposed condensed spectrum is intended to organize condensed parameter
families of microscopic phases.  Similar grading patterns make comparison
reasonable.  They do not make comparison automatic.

Transition charges expose the same issue in a smaller calculation.  The expression
$E^{q+1}(B,B\setminus\Sigma)$ is exact algebraic topology.  Its physical
content begins only after a microscopic phase invariant lands in $E^q(U)$.
Excision localizes the obstruction, and a Thom isomorphism moves it onto the
critical locus under an orientation hypothesis.  None of these steps can
replace the microscopic comparison map.

The program is therefore neither complete nor empty.  Its analytic
foundations include locality bounds, C*-control, gap stability in established
settings, functional calculus, and explicit low-dimensional realizations.
Its categorical assembly is a precise proposal with identifiable theorem
obligations.  That is enough to support a focused research program without
claiming a classification that has not been built.

\section{Conclusion}

The five-stage sequence can be made mathematically coherent if every stage
keeps its hypotheses and every arrow keeps its own proof obligation:
\[
\boxed{
\begin{array}{c}
\text{local quantum interactions}\\
\downarrow\\
\text{condensed moduli stack of Hamiltonians}\\
\downarrow\\
\text{uniformly gapped propositional substack}\\
\downarrow\\
\text{stabilized phase }\infty\text{-groupoid}\\
\downarrow\\
\text{invertible condensed phase spectrum}.
\end{array}}
\]
Locality comes from uniform interaction estimates.  Positivity and norms
come from a declared C*-algebraic setting.  A thermodynamic gap comes from a
quantitative witness.  The qualitative gapped substack is the propositional
image of the witness fibration.  Phases arise only after controlled path
localization and stabilization.  The spectrum arises only after coherent
stacking, restriction to physical units, and internal spectrification.

Under these conditions, $\pi_0$ records invertible stabilized phases,
$\pi_1$ records loops that become pumps when analytic transport is supplied,
and higher homotopy records higher families that become defects only through
a family-to-defect theorem.  The gapless locus is pulled back from the
propositional gapped subobject.  A relative class
$\partial\nu\in E^{q+1}(B,B\setminus\Sigma)$ measures failure of an invariant
to extend, while a physical charge also needs a microscopic comparison.

The proposed spectrum must remain distinct from operator K-theory, Kubota's
microscopic Omega-spectrum, and the Freed-Hopkins field-theory target until
comparison theorems say otherwise.  Explicit models can test those maps.
They cannot be replaced by matching tables of groups.  The next advances
should therefore prove arrows, not rename their endpoints.

\appendix

\section{The five stage data contract}

This appendix records a compact contract suitable for formalization.

\subsection{Local interaction record}

A local interaction family record contains
\begin{equation}
(\Gamma,d,\A_{\mathrm{loc}},\Phi,F,C_F,T_F,G,\alpha),
\end{equation}
where $C_F$ is a common interaction-norm bound and $T_F$ is a common tail
certificate.  Restriction may change the parameter object but may not weaken
the declared bound silently.

\subsection{Hamiltonian moduli record}

A Hamiltonian moduli record adds a test site, object and morphism categories,
restriction functors, descent data, and size conventions.  The object record
alone defines a sheaf of data at best.  Higher moduli require mapping anima
and coherent descent.

\subsection{Gap record}

A gap witness contains
\begin{equation}
(\gamma,L_0,\text{sector},\text{boundary},\text{provider})
\end{equation}
or its GNS analogue.  The qualitative flag is the truncation of the existence
of such a record.  The flag never substitutes for the provider during an
analytic proof.

\subsection{Phase record}

A phase record lists its allowed paths, local circuits, quasi-local
automorphisms, lattice changes, ancillas, and symmetry rules.  Two papers
that use different records may obtain different phase sets without
contradiction.

\subsection{Spectrum record}

A spectrum record lists the symmetric monoidal operation, unit, physical
inverse witness, $E_\infty$ coherences, test-object descent, and spatial
suspension maps.  A sequence of abelian groups is not a substitute for this
record.

\section{Comparison theorem checklist}

Before identifying two phase constructions, a comparison should answer all
of the following questions.

\begin{longtable}{>{\raggedright\arraybackslash}p{0.25\textwidth}
                  >{\raggedright\arraybackslash}p{0.65\textwidth}}
\toprule
Field & Required answer \\
\midrule
\endfirsthead
\toprule
Field & Required answer \\
\midrule
\endhead
Source objects & Local interactions, fields, algebras, or phase objects in the source \\
Source morphisms & Gapped paths, quasi-local maps, field-theory deformations, or stable maps \\
Target objects & Exact target category or spectrum, including grading conventions \\
Parameter site & Smooth manifolds, compact spaces, profinite sets, or condensed objects \\
Analytic domain & Locality norm, gap, regularity, symmetry, and boundary hypotheses \\
Formula & Construction of the map on objects and parameter families \\
Naturality & Compatibility with pullback and restriction \\
Homotopy invariance & Proof that source equivalences map to target equivalences \\
Stacking & Symmetric monoidal compatibility and behavior of units \\
Suspension & Compatibility with dimension shifts and spectrum structure \\
Defects & Treatment of boundaries, higher morphisms, and anomaly inflow \\
Faithfulness & Which distinct microscopic phases can the target merge? \\
Fullness & Which target morphisms lift to microscopic morphisms? \\
Surjectivity & Which target objects have explicit microscopic realizations? \\
Lost data & Energy scales, dispersion, correlation lengths, boundaries, or interactions forgotten by the map \\
Inverse & Explicit construction and proof, if an equivalence is claimed \\
\bottomrule
\end{longtable}

\section{Relative charge bookkeeping}

Let $j:U\hookrightarrow B$.  The relevant part of the exact sequence is
\begin{equation}
\begin{tikzcd}[column sep=large]
E^q(B) \arrow[r,"j^*"]
& E^q(U) \arrow[r,"\partial"]
& E^{q+1}(B,U) \arrow[r]
& E^{q+1}(B).
\end{tikzcd}
\end{equation}
The checks are:
\begin{enumerate}[leftmargin=2.2em]
  \item the connecting map raises degree by one;
  \item exactness gives $\partial j^*=0$;
  \item excision moves the pair to a neighborhood of $\Sigma$;
  \item a rank-$c$ $E$-orientation shifts the localized degree down by $c$;
  \item no orientation means a twisted Thom target;
  \item a physical name requires a microscopic comparison map.
\end{enumerate}

For a codimension-one transition locus, the Thom shift returns degree $q$.
For an isolated codimension-three degeneracy, it gives degree $q-2$ on the
point, equivalently a linking class on an $S^2$ in the normal coordinates.

\section{Status-sensitive glossary}

\begin{longtable}{>{\raggedright\arraybackslash}p{0.23\textwidth}
                  >{\raggedright\arraybackslash}p{0.68\textwidth}}
\toprule
Term & Meaning in this paper \\
\midrule
\endfirsthead
\toprule
Term & Meaning in this paper \\
\midrule
\endhead
Local family & Parameter family with continuous local coefficients and common decay bounds \\
Condensed moduli stack & Proposed anima-valued sheaf of Hamiltonian objects and controlled equivalences \\
Gap witness & Quantitative finite-volume or GNS certificate with a positive common lower bound \\
Witness fibration & Dependent object of Hamiltonians together with chosen certificates \\
Propositional image & Qualitative subobject asserting only that some witness exists \\
Thermodynamic gap & Positive bulk-gap lower bound uniform in the volume convention \\
Phase & Equivalence class in the chosen localized and stabilized microscopic category \\
Invertible phase & Phase with a represented stacking inverse \\
Condensed spectrum & Proposed internal connective spectrum built from the Picard phase object \\
Operator K-class & Stable invariant of a specified operator algebra, not a microscopic reconstruction \\
Kubota spectrum & Established microscopic Omega-spectrum in Kubota's smooth almost-local setting \\
Freed-Hopkins target & Stable-homotopy classification of reflection-positive invertible field theories under its hypotheses \\
Gapless locus & Complement of the pulled-back propositional gapped locus when that complement is represented \\
Relative class & Connecting image $\partial\nu$ in the cohomology of a pair \\
Physical transition charge & Relative class equipped with a microscopic invariant and natural comparison map \\
Realized class & Abstract label with a complete microscopic witness in the stated rules \\
\bottomrule
\end{longtable}

\section{Source papers and division of labor}

The five companion papers supply separate theorem packages.
\begin{enumerate}[leftmargin=2.2em]
  \item \emph{Locality and Lieb-Robinson Bounds for Condensed Families}
  isolates uniform interaction norms, family continuity, thermodynamic
  dynamics, and the moving-defect counterexample.
  \item \emph{Condensed Observable Algebras} treats C*-norms, positivity,
  finite descent, state spaces, crossed products, and Aoki's solidification
  bridge.
  \item \emph{Uniform Thermodynamic Spectral Gaps in Condensed Families}
  separates finite-volume and GNS witnesses and distinguishes the witness
  fibration from its propositional image.
  \item \emph{From Microscopic Lattice Hamiltonians to Effective Field
  Theories} proves the controlled free functional-calculus map, works out the
  SSH transition, and documents information loss and the BDI reduction.
  \item \emph{Physical Realizability of Invertible Phase Classes} compares
  the Freed-Hopkins, Kubota, and proposed condensed targets, verifies three
  low-dimensional witnesses, and states the exact relative-charge
  hypotheses.
\end{enumerate}

The synthesis uses their accepted claims without promoting any companion
proposal to an established theorem.

\begin{thebibliography}{99}

\bibitem{Aoki2024}
K.~Aoki,
\emph{(Semi)topological K-theory via solidification},
arXiv:2409.01462, 2024.

\bibitem{AKLT1987}
I.~Affleck, T.~Kennedy, E.~H.~Lieb, and H.~Tasaki,
\emph{Rigorous results on valence-bond ground states in antiferromagnets},
Physical Review Letters \textbf{59} (1987), 799 to 802,
\url{https://doi.org/10.1103/PhysRevLett.59.799}.

\bibitem{BachmannNachtergaele2014}
S.~Bachmann and B.~Nachtergaele,
\emph{On gapped phases with a continuous symmetry and boundary operators},
Journal of Statistical Physics \textbf{154} (2014), 91 to 112,
\url{https://arxiv.org/abs/1307.0716}.

\bibitem{BMNS2012}
S.~Bachmann, S.~Michalakis, B.~Nachtergaele, and R.~Sims,
\emph{Automorphic equivalence within gapped phases of quantum lattice systems},
Communications in Mathematical Physics \textbf{309} (2012), 835 to 871,
\url{https://arxiv.org/abs/1102.0842}.

\bibitem{Blackadar1998}
B.~Blackadar,
\emph{K-Theory for Operator Algebras}, second edition,
Cambridge University Press, 1998.

\bibitem{Scholze2026}
P.~Scholze,
\emph{Lectures on Condensed Mathematics},
arXiv:2605.03658, 2026.

\bibitem{FidkowskiKitaev2010}
L.~Fidkowski and A.~Kitaev,
\emph{The effects of interactions on the topological classification of free
fermion systems},
Physical Review B \textbf{81} (2010), 134509,
\url{https://doi.org/10.1103/PhysRevB.81.134509}.

\bibitem{FreedHopkins2021}
D.~S.~Freed and M.~J.~Hopkins,
\emph{Reflection positivity and invertible topological phases},
Geometry and Topology \textbf{25} (2021), 1165 to 1330,
\url{https://doi.org/10.2140/gt.2021.25.1165}.

\bibitem{HastingsKoma2006}
M.~B.~Hastings and T.~Koma,
\emph{Spectral gap and exponential decay of correlations},
Communications in Mathematical Physics \textbf{265} (2006), 781 to 804,
\url{https://doi.org/10.1007/s00220-006-0030-4}.

\bibitem{Kitaev2001}
A.~Y.~Kitaev,
\emph{Unpaired Majorana fermions in quantum wires},
Physics-Uspekhi \textbf{44} (2001), 131 to 136,
\url{https://arxiv.org/abs/cond-mat/0010440}.

\bibitem{Kubota2025}
Y.~Kubota,
\emph{Stable homotopy theory of invertible gapped quantum spin systems I:
Kitaev's Omega-spectrum},
arXiv:2503.12618, 2025.

\bibitem{LiebRobinson1972}
E.~H.~Lieb and D.~W.~Robinson,
\emph{The finite group velocity of quantum spin systems},
Communications in Mathematical Physics \textbf{28} (1972), 251 to 257.

\bibitem{NachtergaeleSims2010}
B.~Nachtergaele and R.~Sims,
\emph{Lieb-Robinson bounds in quantum many-body physics},
Contemporary Mathematics \textbf{529} (2010), 141 to 176,
\url{https://arxiv.org/abs/1004.2086}.

\bibitem{ProdanSchulzBaldes2016}
E.~Prodan and H.~Schulz-Baldes,
\emph{Bulk and Boundary Invariants for Complex Topological Insulators},
Springer, 2016.

\bibitem{TeoKane2010}
J.~C.~Y.~Teo and C.~L.~Kane,
\emph{Topological defects and gapless modes in insulators and superconductors},
Physical Review B \textbf{82} (2010), 115120,
\url{https://doi.org/10.1103/PhysRevB.82.115120}.

\bibitem{Thiang2016}
G.~C.~Thiang,
\emph{On the K-theoretic classification of topological phases of matter},
Annales Henri Poincare \textbf{17} (2016), 757 to 794,
\url{https://doi.org/10.1007/s00023-015-0418-9}.

\end{thebibliography}

\end{document}
